\documentclass[12pt]{article}
\usepackage[UTF8]{inputenc}
\usepackage{thaisupp}
\usepackage{enumitem}
\begin{document}
\title{วงจรไฟฟ้ากระแสตรง}
\maketitle
\section*{In-class Questions}
\begin{enumerate}[label=\arabic*.]
\item กระแสไฟฟ้าคือข้อใด? \hfill \textbf{\small (Main)}
\begin{enumerate}[label=)]
\item ก) การเคลื่อนที่ของอิเล็กตรอนอย่างเดียว
\item ข) การเคลื่อนที่ของประจุไฟฟ้าผ่านหน้าตัดใดๆ ต่อหนึ่งหน่วยเวลา
\item ค) ความต่างศักย์ระหว่างสองจุด
\item ง) พลังงานที่ใช้ต่อหน่วยประจุ
\end{enumerate}
\item กระแสไฟฟ้า 2 แอมแปร์ ไหลผ่านตัวนำเป็นเวลา 10 วินาที ประจุที่ไหลผ่านหน้าตัดมีค่าเท่าใด? \hfill \textbf{\small (Main)}
\begin{enumerate}[label=)]
\item ก) 0.2 C
\item ข) 2 C
\item ค) 5 C
\item ง) 20 C
\end{enumerate}
\item กฎของโอห์ม (Ohm's Law) คือข้อใด? \hfill \textbf{\small (Main)}
\begin{enumerate}[label=)]
\item ก) V = IR
\item ข) V = I/R
\item ค) V = R/I
\item ง) V = I²R
\end{enumerate}
\item ผ่านความต้านทาน 100 Ω ด้วยกระแส 0.5 A ความต่างศักย์คร่อมตัวต้านทานเป็นเท่าใด? \hfill \textbf{\small (Main)}
\begin{enumerate}[label=)]
\item ก) 50 V
\item ข) 200 V
\item ค) 0.005 V
\item ง) 5 V
\end{enumerate}
\item ตัวต้านทานสองตัว 4 Ω และ 6 Ω ต่ออนุกรมกัน ความต้านทานรวมเป็นเท่าใด? \hfill \textbf{\small (Main)}
\begin{enumerate}[label=)]
\item ก) 2.4 Ω
\item ข) 4 Ω
\item ค) 6 Ω
\item ง) 10 Ω
\end{enumerate}
\item ตัวต้านทานสองตัว 4 Ω และ 6 Ω ต่อขนานกัน ความต้านทานรวมเป็นเท่าใด? \hfill \textbf{\small (Main)}
\begin{enumerate}[label=)]
\item ก) 2.4 Ω
\item ข) 4 Ω
\item ค) 6 Ω
\item ง) 10 Ω
\end{enumerate}
\item กำลังไฟฟ้า P บนตัวต้านทาน R ที่มีกระแส I ไหลผ่าน มีค่าเท่าใด? \hfill \textbf{\small (Main)}
\begin{enumerate}[label=)]
\item ก) P = IV
\item ข) P = I²R
\item ค) P = V²/R
\item ง) ถูกทุกข้อ
\end{enumerate}
\item หลอดไฟขนาด 100 W ใช้ที่ความต่างศักย์ 220 V ความต้านทานของหลอดไฟเป็นเท่าใด? \hfill \textbf{\small (Main)}
\begin{enumerate}[label=)]
\item ก) 2.2 Ω
\item ข) 22 Ω
\item ค) 220 Ω
\item ง) 484 Ω
\end{enumerate}
\item กฎกระแสของเคอร์ชฮอฟฟ์ (Kirchhoff's Current Law) ระบุว่าผลรวมพีชคณิตของกระแสที่จุดรวม (junction) เป็นเท่าใด? \hfill \textbf{\small (Main)}
\begin{enumerate}[label=)]
\item ก) เท่ากับผลรวมของแรงเคลื่อน
\item ข) เท่ากับศูนย์
\item ค) เท่ากับกระแสรวมที่ไหลออก
\item ง) เท่ากับความต้านทานรวม
\end{enumerate}
\item กฎแรงเคลื่อนของเคอร์ชฮอฟฟ์ (Kirchhoff's Voltage Law) ระบุว่าผลรวมพีชคณิตของแรงเคลื่อนไฟฟ้ารอบห่วง (loop) เป็นเท่าใด? \hfill \textbf{\small (Main)}
\begin{enumerate}[label=)]
\item ก) เท่ากับกระแส
\item ข) เท่ากับศูนย์
\item ค) เท่ากับผลรวมของความต้านทาน
\item ง) เท่ากับกำลัง
\end{enumerate}
\item ตัวต้านทาน 3 ตัว 6 Ω ต่ออนุกรมกัน ความต้านทานรวมเป็นเท่าใด? \hfill \textbf{\small (Main)}
\begin{enumerate}[label=)]
\item ก) 2 Ω
\item ข) 6 Ω
\item ค) 9 Ω
\item ง) 18 Ω
\end{enumerate}
\item ตัวต้านทาน 3 ตัว 6 Ω ต่อขนานกัน ความต้านทานรวมเป็นเท่าใด? \hfill \textbf{\small (Main)}
\begin{enumerate}[label=)]
\item ก) 2 Ω
\item ข) 6 Ω
\item ค) 9 Ω
\item ง) 18 Ω
\end{enumerate}
\item ถ้า 6×10¹⁸ อิเล็กตรอนไหลผ่านจุดใดในเวลา 2 วินาที กระแสมีค่าเท่าใด? (e = 1.6×10⁻¹⁹ C) \hfill \textbf{\small (Main)}
\begin{enumerate}[label=)]
\item ก) 0.48 A
\item ข) 0.96 A
\item ค) 4.8 A
\item ง) 9.6 A
\end{enumerate}
\item ตัวต้านทาน 20 Ω ผ่านกระแส 3 A เป็นเวลา 1 นาที พลังงานความร้อนที่เกิดขึ้นเป็นเท่าใด? \hfill \textbf{\small (Main)}
\begin{enumerate}[label=)]
\item ก) 60 J
\item ข) 180 J
\item ค) 3600 J
\item ง) 10800 J
\end{enumerate}
\item ลวดตัวนำยาว 10 m มีพื้นที่หน้าตัด 2 mm² มีความต้านทาน 5 Ω ความต้านทานของลวดชนิดเดียวกันยาว 25 m และพื้นที่หน้าตัด 5 mm² เป็นเท่าใด? \hfill \textbf{\small (Main)}
\begin{enumerate}[label=)]
\item ก) 2.5 Ω
\item ข) 5 Ω
\item ค) 10 Ω
\item ง) 25 Ω
\end{enumerate}
\item ตัวต้านทาน R₁ = 2 Ω, R₂ = 3 Ω, R₃ = 6 Ω ต่อแบบผสม โดย R₁ อนุกรมกับกลุ่ม R₂ ขนาน R₃ ความต้านทานรวมเป็นเท่าใด? \hfill \textbf{\small (Main)}
\begin{enumerate}[label=)]
\item ก) 2 Ω
\item ข) 3 Ω
\item ค) 4 Ω
\item ง) 5 Ω
\end{enumerate}
\item วงจรประกอบด้วยแบตเตอรี่ 12 V ต่ออนุกรมกับตัวต้านทาน 4 Ω และ 8 Ω กระแสที่ไหลในวงจรเป็นเท่าใด? \hfill \textbf{\small (Main)}
\begin{enumerate}[label=)]
\item ก) 0.5 A
\item ข) 1 A
\item ค) 1.5 A
\item ง) 2 A
\end{enumerate}
\item เครื่องทำความร้อนขนาด 1000 W ใช้กับไฟบ้าน 220 V กระแสที่ใช้มีค่าเท่าใด? \hfill \textbf{\small (Main)}
\begin{enumerate}[label=)]
\item ก) 0.22 A
\item ข) 2.27 A
\item ค) 4.55 A
\item ง) 22 A
\end{enumerate}
\item ตัวต้านทาน 4 Ω ต่ออนุกรมกับกลุ่มตัวต้านทาน 2 Ω และ 3 Ω ที่ต่อขนานกัน ความต้านทานรวมเป็นเท่าใด? \hfill \textbf{\small (Main)}
\begin{enumerate}[label=)]
\item ก) 5.2 Ω
\item ข) 6 Ω
\item ค) 7.2 Ω
\item ง) 9 Ω
\end{enumerate}
\item วงจรประกอบด้วยแบตเตอรี่ EMF 24 V ความต้านภายใน 2 Ω ต่อกับตัวต้านทานภายนอก 6 Ω กระแสที่ไหลและความต่างศักย์คร่อมขั้วแบตเตอรี่เป็นเท่าใด? \hfill \textbf{\small (Main)}
\begin{enumerate}[label=)]
\item ก) I = 3 A, V = 18 V
\item ข) I = 3 A, V = 24 V
\item ค) I = 4 A, V = 18 V
\item ง) I = 4 A, V = 24 V
\end{enumerate}
\item สภาพนำไฟฟ้า (Conductivity) σ ของตัวนำมีความสัมพันธ์กับสภาพต้านทานไฟฟ้า (Resistivity) ρ อย่างไร? \hfill \textbf{\small (Main)}
\begin{enumerate}[label=)]
\item ก) σ = 1/ρ
\item ข) σ = ρ
\item ค) σ = ρ²
\item ง) σ = 1/ρ²
\end{enumerate}
\item หลอดไฟ 60 W ต่อกับแบตเตอรี่ 12 V กระแสที่ไหลผ่านหลอดไฟมีค่าเท่าใด? \hfill \textbf{\small (Main)}
\begin{enumerate}[label=)]
\item ก) 0.2 A
\item ข) 5 A
\item ค) 72 A
\item ง) 720 A
\end{enumerate}
\item วงจรมีตัวต้านทาน 2 Ω, 3 Ω, 5 Ω ต่อขนานกัน กับแบตเตอรี่ 15 V กระแสรวมที่ไหลจากแบตเตอรี่เป็นเท่าใด? \hfill \textbf{\small (Main)}
\begin{enumerate}[label=)]
\item ก) 1.5 A
\item ข) 4 A
\item ค) 5 A
\item ง) 15 A
\end{enumerate}
\item วงจรมีแบตเตอรี่ 2 วงจร วงที่ 1 มี EMF 10 V และความต้านภายใน 1 Ω วงที่ 2 มี EMF 5 V และความต้านภายใน 1 Ω ต่อขนานกัน ต่อกับตัวต้านทานภายนอก 4 Ω EMF รวมเป็นเท่าใด? \hfill \textbf{\small (Main)}
\begin{enumerate}[label=)]
\item ก) 5 V
\item ข) 7.5 V
\item ค) 10 V
\item ง) 15 V
\end{enumerate}
\item เครื่องทำความร้อน 2 เครื่อง กำลัง 500 W ที่ 220 V ต่ออนุกรมกัน แล้วต่อกับไฟ 220 V กำลังที่ใช้จริงเป็นเท่าใด? \hfill \textbf{\small (Main)}
\begin{enumerate}[label=)]
\item ก) 125 W
\item ข) 250 W
\item ค) 500 W
\item ง) 1000 W
\end{enumerate}
\item ทิศทางของกระแสไฟฟ้าตามข้อตกลง (Conventional Current) คือทิศทางใด? \hfill \textbf{\small (Extra)}
\begin{enumerate}[label=)]
\item ก) ทิศทางที่อิเล็กตรอนเคลื่อนที่
\item ข) ทิศทางที่ประจุบวกเคลื่อนที่ หรือตรงข้ามทิศทางที่อิเล็กตรอนเคลื่อนที่
\item ค) ทิศทางสนามไฟฟ้า
\item ง) ทิศทางแรงโน้มถ่วง
\end{enumerate}
\item หน่วยของกระแสไฟฟ้าในหน่วย SI คือข้อใด? \hfill \textbf{\small (Extra)}
\begin{enumerate}[label=)]
\item ก) โวลต์ (V)
\item ข) แอมแปร์ (A)
\item ค) โอห์ม (Ω)
\item ง) วัตต์ (W)
\end{enumerate}
\item ตัวต้านทานมีแถบสี แดง แดง น้ำตาล ค่าความต้านทานเป็นเท่าใด? \hfill \textbf{\small (Extra)}
\begin{enumerate}[label=)]
\item ก) 22 Ω
\item ข) 220 Ω
\item ค) 2.2 kΩ
\item ง) 22 kΩ
\end{enumerate}
\item สารในข้อใดเป็นตัวนำไฟฟ้าที่ดีที่สุด? \hfill \textbf{\small (Extra)}
\begin{enumerate}[label=)]
\item ก) ทองแดง (Cu)
\item ข) อะลูมิเนียม (Al)
\item ค) เงิน (Ag)
\item ง) ทองคำ (Au)
\end{enumerate}
\item การต่อตัวต้านทานแบบใดที่กระแสในวงจรทุกตัวเท่ากัน? \hfill \textbf{\small (Extra)}
\begin{enumerate}[label=)]
\item ก) การต่อขนาน
\item ข) การต่ออนุกรม
\item ค) การต่อผสม
\item ง) การต่อแบบbridge
\end{enumerate}
\item การต่อตัวต้านทานแบบใดที่ความต่างศักย์คร่อมตัวต้านทานทุกตัวเท่ากัน? \hfill \textbf{\small (Extra)}
\begin{enumerate}[label=)]
\item ก) การต่ออนุกรม
\item ข) การต่อขนาน
\item ค) การต่อผสม
\item ง) การต่อแบบbridge
\end{enumerate}
\item พลังงานไฟฟ้า 1 หน่วย (หน่วยการไฟฟ้า) เท่ากับกี่จูล? \hfill \textbf{\small (Extra)}
\begin{enumerate}[label=)]
\item ก) 3.6×10³ J
\item ข) 3.6×10⁶ J
\item ค) 1 J
\item ง) 10⁶ J
\end{enumerate}
\item เครื่องใช้ไฟฟ้า 500 W ใช้งาน 8 ชั่วโมง กี่หน่วยของไฟฟ้า? \hfill \textbf{\small (Extra)}
\begin{enumerate}[label=)]
\item ก) 0.0625 หน่วย
\item ข) 0.5 หน่วย
\item ค) 4 หน่วย
\item ง) 16 หน่วย
\end{enumerate}
\item แบตเตอรี่ EMF 12 V ความต้านภายใน 0.5 Ω �ต่อกับวงจร กระแสลัดวงจร (short circuit current) เป็นเท่าใด? \hfill \textbf{\small (Extra)}
\begin{enumerate}[label=)]
\item ก) 0.5 A
\item ข) 6 A
\item ค) 12 A
\item ง) 24 A
\end{enumerate}
\item ลวดทองแดงเส้นหนึ่งยาว 100 m มีพื้นที่หน้าตัด 1 mm² มีความต้านทานเท่าใด? (ρ\_Cu = 1.68×10⁻⁸ Ω·m) \hfill \textbf{\small (Extra)}
\begin{enumerate}[label=)]
\item ก) 0.168 Ω
\item ข) 1.68 Ω
\item ค) 16.8 Ω
\item ง) 168 Ω
\end{enumerate}
\item ตัวต้านทาน 2 Ω, 3 Ω, 6 Ω ต่ออนุกรมกัน แล้วต่อขนานกับตัวต้านทาน 1 Ω ความต้านทานรวมเป็นเท่าใด? \hfill \textbf{\small (Extra)}
\begin{enumerate}[label=)]
\item ก) 1 Ω
\item ข) 2 Ω
\item ค) 6 Ω
\item ง) 11 Ω
\end{enumerate}
\item แบตเตอรี่ 6 V ความต้านภายใน 1 Ω จ่ายกระแสให้ตัวต้านทานภายนอก R ถ้ากำลังที่ใช้ในตัวต้านทานภายนอกมากที่สุด ค่า R ต้องเป็นเท่าใด? \hfill \textbf{\small (Extra)}
\begin{enumerate}[label=)]
\item ก) 0.5 Ω
\item ข) 1 Ω
\item ค) 2 Ω
\item ง) 6 Ω
\end{enumerate}
\item หลอดไฟ A ขนาด 100 W ที่ 220 V และหลอดไฟ B ขนาด 60 W ที่ 220 V ต่อขนานกันกับไฟ 220 V ข้อใดถูกต้อง? \hfill \textbf{\small (Extra)}
\begin{enumerate}[label=)]
\item ก) หลอด A สว่างกว่า
\item ข) หลอด B สว่างกว่า
\item ค) หลอดทั้งสองสว่างเท่ากัน
\item ง) หลอดไฟดับทั้งคู่
\end{enumerate}
\item วงจรมีตัวต้านทาน 6 ตัว 12 Ω ต่อเป็นรูปลูกบาศก์ (cube) ที่ขอบทุกขอบ ความต้านทานระหว่างจุดยอดที่เป็นมุมตรงข้ามเป็นเท่าใด? \hfill \textbf{\small (Extra)}
\begin{enumerate}[label=)]
\item ก) 2 Ω
\item ข) 4 Ω
\item ค) 6 Ω
\item ง) 8 Ω
\end{enumerate}
\item อิเล็กตรอนเคลื่อนที่ในลวดตัวนำด้วยความเร็วลอยเลื่อน (drift velocity) v เมื่อเพิ่มความต่างศักย์เป็น 2 เท่า ความเร็วลอยเลื่อนจะเป็นเท่าใด? \hfill \textbf{\small (Extra)}
\begin{enumerate}[label=)]
\item ก) v/2
\item ข) v
\item ค) 2v
\item ง) 4v
\end{enumerate}
\item วงจรประกอบด้วยแบตเตอรี่ 2 วงจร วงที่ 1: ε₁ = 6 V, r₁ = 1 Ω วงที่ 2: ε₂ = 12 V, r₂ = 1 Ω ต่ออนุกรมกันช่วยกัน (รวมแรงเคลื่อน) ต่อกับ R = 5 Ω กระแสที่ไหลเป็นเท่าใด? \hfill \textbf{\small (Extra)}
\begin{enumerate}[label=)]
\item ก) 1 A
\item ข) 2 A
\item ค) 3 A
\item ง) 4 A
\end{enumerate}
\item แบตเตอรี่ 12 V ความต้านภายใน 2 Ω จ่ายกระแสให้วงจร ถ้ากำลังที่จ่ายให้ภายนอกเป็น 10 W กระแสที่ไหลเป็นเท่าใด? \hfill \textbf{\small (Extra)}
\begin{enumerate}[label=)]
\item ก) 1 A
\item ข) 2 A
\item ค) 3 A
\item ง) 4 A
\end{enumerate}
\item สารกึ่งตัวนำ (Semiconductor) มีค่าสภาพต้านทานไฟฟ้าอยู่ในช่วงใด? \hfill \textbf{\small (Extra)}
\begin{enumerate}[label=)]
\item ก) 10⁻⁸ ถึง 10⁻⁶ Ω·m
\item ข) 10⁻⁶ ถึง 10⁻⁴ Ω·m
\item ค) 10⁻⁵ ถึง 10⁶ Ω·m
\item ง) 10⁻⁴ ถึง 10⁶ Ω·m
\end{enumerate}
\item ตัวต้านทาน 5 ตัว 10 Ω ต่อขนานกัน ความต้านทานรวมเป็นเท่าใด? \hfill \textbf{\small (Extra)}
\begin{enumerate}[label=)]
\item ก) 0.5 Ω
\item ข) 2 Ω
\item ค) 10 Ω
\item ง) 50 Ω
\end{enumerate}
\item วงจรมีแบตเตอรี่ ε = 24 V, r = 2 Ω และตัวต้านทานภายนอก 2 ตัว 4 Ω กับ 12 Ω ต่อขนานกัน กระแสที่ไหลผ่านตัวต้านทาน 4 Ω เป็นเท่าใด? \hfill \textbf{\small (Extra)}
\begin{enumerate}[label=)]
\item ก) 1.5 A
\item ข) 2 A
\item ค) 3 A
\item ง) 4 A
\end{enumerate}
\item มอเตอร์ไฟฟ้าขนาด 1 HP (1 HP ≈ 746 W) มีประสิทธิภาพ 80\% ต่อกับไฟ 220 V กระแสที่ใช้มีค่าเท่าใด? \hfill \textbf{\small (Extra)}
\begin{enumerate}[label=)]
\item ก) 2.55 A
\item ข) 3.39 A
\item ค) 4.24 A
\item ง) 5.09 A
\end{enumerate}
\item ตัวต้านทาน 2 ตัวต่ออนุกรมใช้กำลังรวม 18 W กับแบตเตอรี่ 9 V ถ้าตัวต้านทานตัวหนึ่งเป็น 2 Ω อีกตัวมีค่าเท่าใด? \hfill \textbf{\small (Extra)}
\begin{enumerate}[label=)]
\item ก) 1 Ω
\item ข) 2 Ω
\item ค) 3 Ω
\item ง) 4 Ω
\end{enumerate}
\item กระแสไฟฟ้าที่ไหลในตัวนำโลหะเกิดจากการเคลื่อนที่ของอะไร? \hfill \textbf{\small (Extra)}
\begin{enumerate}[label=)]
\item ก) โปรตอน
\item ข) นิวตรอน
\item ค) อิเล็กตรอนอิสระ
\item ง) ไอออน
\end{enumerate}
\item แบตเตอรี่ 24 V ความต้านภายใน 2 Ω จ่ายกระแสให้วงจรที่มีตัวต้านทานหลายตัว ถ้ากำลังที่สูญเสียในแบตเตอรี่เป็น 2 เท่าของกำลังที่จ่ายให้ภายนอก กระแสที่ไหลเป็นเท่าใด? \hfill \textbf{\small (Extra)}
\begin{enumerate}[label=)]
\item ก) 2 A
\item ข) 3 A
\item ค) 4 A
\item ง) 6 A
\end{enumerate}
\item แถบสีของตัวต้านทาน: น้ำตาล เขียว แดง ทอง ค่าความต้านทานเป็นเท่าใด? \hfill \textbf{\small (Extra)}
\begin{enumerate}[label=)]
\item ก) 15×10² ±5\% Ω
\item ข) 15×10² ±10\% Ω
\item ค) 15×10³ ±5\% Ω
\item ง) 1.5×10³ ±5\% Ω
\end{enumerate}
\item ตัวต้านทาน 3 ตัว R ต่อขนานกัน ความต้านทานรวมเป็นเท่าใด? \hfill \textbf{\small (Extra)}
\begin{enumerate}[label=)]
\item ก) R/3
\item ข) R
\item ค) 3R
\item ง) R/9
\end{enumerate}
\item หลอดไฟ 40 W ที่ 220 V กับหลอดไฟ 60 W ที่ 220 V ต่ออนุกรมกันกับไฟ 220 V ข้อใดถูกต้อง? \hfill \textbf{\small (Extra)}
\begin{enumerate}[label=)]
\item ก) หลอด 40 W สว่างกว่า
\item ข) หลอด 60 W สว่างกว่า
\item ค) หลอดทั้งสองสว่างเท่ากัน
\item ง) หลอดไฟดับทั้งคู่
\end{enumerate}
\item วงจรมีแหล่งจ่ายไฟ 18 V ต่อกับ R₁ = 3 Ω, R₂ = 6 Ω อนุกรมกัน และ R₃ = 4 Ω ต่อขนานกับกลุ่ม R₁-R₂ ความต่างศักย์คร่อม R₃ เป็นเท่าใด? \hfill \textbf{\small (Extra)}
\begin{enumerate}[label=)]
\item ก) 4 V
\item ข) 6 V
\item ค) 9 V
\item ง) 12 V
\end{enumerate}
\item สภาพต้านทานไฟฟ้าของโลหะมีค่าขึ้นอยู่กับปัจจัยใด? \hfill \textbf{\small (Extra)}
\begin{enumerate}[label=)]
\item ก) อุณหภูมิเท่านั้น
\item ข) พื้นที่หน้าตัดเท่านั้น
\item ค) อุณหภูมิและชนิดของโลหะ
\item ง) ความยาวเท่านั้น
\end{enumerate}
\item แบตเตอรี่ 10 V ความต้านภายใน 1 Ω จ่ายกระแสให้ตัวต้านทาน R กำลังที่ใช้ใน R มากที่สุดเท่าใด? \hfill \textbf{\small (Extra)}
\begin{enumerate}[label=)]
\item ก) 10 W
\item ข) 20 W
\item ค) 25 W
\item ง) 50 W
\end{enumerate}
\end{enumerate}
\end{document}
